% P15-U3: exact assumptions for the computation methods in Kass--Raftery Sec. 4.
\chapter{Computing Bayes Factors}
\label{sec:p15-bayes-factor-computation}

A marginal likelihood is an integral, and its numerical approximation is
a further operation with its own error conditions. This chapter develops
the Laplace, one-step Newton, Schwarz/BIC, direct Monte Carlo, importance,
harmonic-mean, mixture, bridge, Laplace--Metropolis, Savage--Dickey and
latent-data methods described by Kass and Raftery
\cite[Sec.~4, pp.~777--781]{KassRaftery1995BayesFactors}.
The model and data-measure conventions are those of
\cref{def:p15-model-comparison-experiment}. Likelihood normalizing constants
remain part of every expression unless a common cancellation is proved.

The sample-information index \(n\), the parameter dimension \(d\), and the
simulation count \(M\) have different roles. Increasing \(M\) does not repair
an invalid large-\(n\) approximation. Increasing \(n\) does not make an
arbitrary simulation scheme converge to its intended distribution.

\section{A Laplace theorem with explicit local and tail conditions}
\label{sec:p15-laplace-hypotheses}

\begin{hbdefinition}{def:p15-laplace-family}{A uniformly controlled interior mode}
Fix an integer \(d\geq1\). For each \(n\geq1\), let
\(\Theta_n\subseteq\mathbb R^d\) be measurable and contain the ball
\(B(\theta_n,\rho)\), for a fixed \(\rho>0\). Consider
\[
 I_n=\int_{\Theta_n}e^{nh_n(\theta)}a_n(\theta)\,d\theta .
\]
Here \(h_n\) is measurable with values in \(\mathbb R\cup\{-\infty\}\),
and \(a_n\) is finite, nonnegative and measurable. On the stated ball,
\(h_n\) is real and \(C^4\), \(a_n\) is \(C^2\), and the following
constants are independent of \(n\):
\begin{enumerate}
\item \(\nabla h_n(\theta_n)=0\), and
\(J_n=-D^2h_n(\theta_n)\) has all eigenvalues in \([c,C]\), with
\(0<c\leq C<\infty\).
\item The operator norms of \(D^3h_n,D^4h_n,Da_n,D^2a_n\) on the ball
are at most \(A<\infty\), and
\(0<a_-\leq a_n(\theta_n)\leq a_+<\infty\).
\item For every \(v\) with \(\|v\|<\rho\),
\(h_n(\theta_n+v)-h_n(\theta_n)\leq-c_0\|v\|^2\), where \(c_0>0\).
\item For constants \(A_T,\eta>0\),
\[
 \int_{\Theta_n\setminus B(\theta_n,\rho)}
 e^{n(h_n(\theta)-h_n(\theta_n))}a_n(\theta)\,d\theta
 \leq A_Te^{-\eta n}.
\]
\end{enumerate}
The norm is Euclidean, derivative norms are its induced multilinear norms,
and \(d\theta\) is \(d\)-dimensional Lebesgue measure. The last condition
controls all mass outside the local mode, including other modes and tails.
These are deterministic conditions on the sequence being approximated.
\end{hbdefinition}

\begin{hbtheorem}{thm:p15-laplace-relative-error}{proved=annex:p15.laplace-relative-error}{Relative Laplace accuracy under explicit hypotheses}
Under \cref{def:p15-laplace-family},
\[
 \widehat I_n=e^{nh_n(\theta_n)}a_n(\theta_n)
                  (2\pi/n)^{d/2}(\det J_n)^{-1/2}
\]
satisfies \(I_n/\widehat I_n=1+O(n^{-1})\).
The implied constant depends only on the displayed uniform bounds and
fixed dimension, not on \(n\). In particular \(I_n\) is positive and finite.
For two such marginal integrals,
\(B_{10}=I_{1,n}/I_{0,n}\) and
\(\widehat B_{10}=\widehat I_{1,n}/\widehat I_{0,n}\) satisfy
\(B_{10}/\widehat B_{10}=1+O(n^{-1})\) and
\(\log B_{10}-\log\widehat B_{10}=O(n^{-1})\).
\end{hbtheorem}

This theorem does not claim the hypotheses for every statistical model or
for dimensions increasing with \(n\). To apply it to random data, one must
establish the stated bounds along the data sequences in question. A label
such as regular model is not a substitute for the local and tail premises.

\section{Posterior-mode and likelihood-mode constructions}
\label{sec:p15-laplace-mode-variants}

For an observed data set write \(L_n(\theta)\) for its full likelihood,
\(\pi(\theta)\) for a proper prior density, and
\(I_n=\int L_n(\theta)\pi(\theta)d\theta\). Logs below are taken only in
the positive local neighborhood of the specified mode.

\begin{hbcorollary}{thm:p15-laplace-mode-variants}{proved=annex:p15.laplace-mode-variants}{Two expansions of the same integral}
For the posterior-mode construction let
\[
 q_n=\log(L_n\pi),\quad \widetilde\theta_n\in\arg\max q_n,\quad
 H_n=-D^2q_n(\widetilde\theta_n).
\]
If \(h_n=q_n/n,a_n=1,\theta_n=\widetilde\theta_n\) satisfy the explicit
Laplace hypotheses, then
\[
 \widehat I_{\rm post}
 =(2\pi)^{d/2}(\det H_n)^{-1/2}
                L_n(\widetilde\theta_n)\pi(\widetilde\theta_n)
\]
has relative error \(O(n^{-1})\).

For the likelihood-mode construction let
\[
 \ell_n=\log L_n,\quad \widehat\theta_n\in\arg\max\ell_n,\quad
 F_n=-D^2\ell_n(\widehat\theta_n).
\]
If \(h_n=\ell_n/n,a_n=\pi,\theta_n=\widehat\theta_n\) satisfy those
hypotheses, then
\[
 \widehat I_{\rm MLE}
 =(2\pi)^{d/2}(\det F_n)^{-1/2}
                  L_n(\widehat\theta_n)\pi(\widehat\theta_n)
\]
also has relative error \(O(n^{-1})\). These are separate applications
with separate hypothesis checks.
\end{hbcorollary}

The first expansion absorbs the prior into the exponent and uses posterior
curvature; the second leaves it as an amplitude and uses likelihood
curvature. A strongly varying prior can make the second less accurate at a
given finite sample size \cite[Secs.~4.1.1--4.1.2]{KassRaftery1995BayesFactors}.

\begin{hblemma}{thm:p15-laplace-curvature-replacement}{proved=annex:p15.laplace-curvature-replacement}{Replacing an information matrix}
Suppose \(F_n=nJ_n\), where \(J_n\) is symmetric with all eigenvalues
at least \(c>0\).
Replace it by \(\widetilde F_n=n\widetilde J_n\), where
\(\widetilde J_n\) is symmetric and
\(\|\widetilde J_n-J_n\|\leq\delta_n<c/2\). Holding the mode, likelihood
and prior ordinate fixed, the changed approximation satisfies
\[
 \left|\log\frac{\widehat I(\widetilde F_n)}
                         {\widehat I(F_n)}\right|
 \leq d\delta_n/c.
\]
Thus a bound \(\delta_n=O(n^{-1/2})\) gives relative error
\(O(n^{-1/2})\), and \(O(n^{-1})\) preserves that smaller error order.
An expected-information matrix may be substituted under such a proved
comparison; the name expected information alone does not supply the bound.
\end{hblemma}

Kass and Raftery state that the expected-information version of
approximation (5) has the larger relative error \(O(n^{-1/2})\)
\cite[Sec.~4.1.2, p.~778]{KassRaftery1995BayesFactors}. The next result
computes a data sequence on which that order is attained exactly, and a
second sequence on which no comparison of the required size holds.

\begin{hbtheorem}{thm:p15-w02b-p15-056-information-replacement-boundary}{proved=annex:p15-w02b-p15-056-infobound}{Expected information at exact order \(n^{-1/2}\), and a failed comparison}
Let \(f(u)=e^{-u}/(1+e^{-u})^2\) be the logistic density and take the
location model \(x\mid\theta\sim f(x-\theta)\) with a
\(\operatorname{normal}(0,1)\) prior on \(\theta\in\mathbb R\). The
expected information per observation is \(1/3\) at every \(\theta\).
For even \(n=2m\), consider a data set with \(m\) observations at
\(a\) and \(m\) at \(-a\), where \(a\geq0\) may depend on \(n\).
\begin{enumerate}
\item The likelihood mode is \(\widehat\theta_n=0\) and the observed
information per observation is \(J_n=2f(a)\).
\item For even \(n\geq38\), let \(a=a_n\) solve
\(2f(a_n)=1/3+n^{-1/2}\). Then \(J_n=1/3+n^{-1/2}\),
\(\widetilde J_n=1/3\), \(\delta_n=n^{-1/2}<c/2\) with \(c=1/3\), the
MLE Laplace hypotheses hold, and
\[
 \log\frac{\widehat I(\widetilde F_n)}{I_n}
 =\tfrac12\log\bigl(1+3n^{-1/2}\bigr)+O(n^{-1}),
 \qquad
 \sqrt n\left(\frac{\widehat I(\widetilde F_n)}{I_n}-1\right)\to\frac32 .
\]
The expected-information approximation therefore has relative error of
exact order \(n^{-1/2}\), while the observed-information approximation
has relative error \(O(n^{-1})\).
\item For every even \(n\), let \(a=\log(11+\sqrt{120})\), so that
\(2f(a)=1/12\). Then \(J_n=1/12\) and \(\|\widetilde J_n-J_n\|=1/4\).
Every admissible eigenvalue bound satisfies \(c\leq1/12\), so the
comparison \(\delta_n<c/2\) of \cref{thm:p15-laplace-curvature-replacement}
never holds. The MLE Laplace hypotheses hold with observed information,
and \(\widehat I(\widetilde F_n)/I_n\to1/2\): the expected-information
approximation has relative error tending to \(-1/2\).
\end{enumerate}
\end{hbtheorem}

The data in the third item are possible observations under the model.
The point is that a replacement of observed by expected information is
justified by a comparison along the data sequence actually observed. The
model's name does not provide it.

\begin{hbmethod}{def:p15-w02b-p15-056-one-step-newton}{The one-step Newton approximation}
For models \(H_k\), \(k=0,1\), with parameters \(\theta_k\) of dimension
\(d_k\), use the notation of Kass and Raftery
\cite[Sec.~4.1.2, eqs.~(7)--(8), p.~778]{KassRaftery1995BayesFactors}.
\(\widehat\theta_k\) is the maximum-likelihood estimate, and
\(F_k=\widehat\Sigma_k^{-1}\) is the observed information at it.
\(\lambda_k(\theta_k)=\log\pi(\theta_k\mid H_k)\) is the log prior
density, and \(\lambda_k'(\theta_k)\) is its gradient. The prior
covariance \(W_k=\operatorname{var}[\theta_k\mid H_k]\) is assumed
finite and positive definite, and \(G_k=W_k^{-1}\). The steps are as
follows.
\begin{enumerate}
\item Compute \(\widehat\theta_k\), \(F_k\), the maximized
log-likelihood \(\ell_k(\widehat\theta_k)\),
\(\lambda_k(\widehat\theta_k)\), \(\lambda_k'(\widehat\theta_k)\) and \(G_k\).
\item Take one Newton step for the log posterior kernel
\(\ell_k+\lambda_k\) from \(\widehat\theta_k\), with curvature matrix
\(-(F_k+G_k)\):
\(\bar\theta_k=\widehat\theta_k+(F_k+G_k)^{-1}\lambda_k'(\widehat\theta_k)\).
\item Form
\[
\begin{split}
 E_k={}&2\lambda_k(\widehat\theta_k)
 +\lambda_k'(\widehat\theta_k)^T(F_k+G_k)^{-1}
   \{2-F_k(F_k+G_k)^{-1}\}\lambda_k'(\widehat\theta_k)\\
 &-\log|F_k+G_k|+d_k\log(2\pi),
\end{split}
\]
where \(2\) inside the braces is twice the identity matrix.
\item With \(\Lambda=2\{\ell_1(\widehat\theta_1)-\ell_0(\widehat\theta_0)\}\),
report \(2\log B_{10}\approx\Lambda+(E_1-E_0)\).
\end{enumerate}
Only quantities at the maximum-likelihood estimate and the prior's first
two moments enter. A parameter-free model uses \(E_k=0\).
\end{hbmethod}

\begin{hbtheorem}{thm:p15-w02b-p15-056-one-step-newton}{proved=annex:p15-w02b-p15-056-newton}{What the one-step formula computes, and its error}
In the setting of \cref{def:p15-w02b-p15-056-one-step-newton}, drop the
index \(k\) and write \(\ell=\ell_n\), \(F=F_n\).
\begin{enumerate}
\item Let \(\ell^{(2)}(\theta)=\ell(\widehat\theta)
-\tfrac12(\theta-\widehat\theta)^TF(\theta-\widehat\theta)\) and
\(\lambda^{(1)}(\theta)=\lambda(\widehat\theta)
+\lambda'(\widehat\theta)^T(\theta-\widehat\theta)\). Then
\[
 2\ell(\widehat\theta)+E
 =2\{\ell^{(2)}(\bar\theta)+\lambda^{(1)}(\bar\theta)\}
 +d\log(2\pi)-\log|F+G| .
\]
Thus \(E\) is twice the logarithm of the posterior-mode formula (4) after
three replacements: the mode by \(\bar\theta\), the posterior curvature by
\(F+G\), and the log likelihood and log prior by their quadratic and
linear expansions at \(\widehat\theta\).
\item Suppose the MLE construction of \cref{thm:p15-laplace-mode-variants}
satisfies \cref{def:p15-laplace-family}, with \(h_n=\ell_n/n\),
\(a_n=\pi\), \(F_n=nJ_n\), and suppose \(\|G\|\leq g\) for a constant
\(g\) independent of \(n\). Then
\[
 2\log I_n=2\ell_n(\widehat\theta_n)+E+O(n^{-1}).
\]
If both models satisfy these conditions,
\(2\log B_{10}=\Lambda+(E_1-E_0)+O(n^{-1})\).
\item In the normal model of \cref{thm:p15-laplace-normal-example}, put
\(D_n=x_n-\mu\) and \(s^2=\sigma^2/n\). Then
\[
 \ell_n(\widehat\theta_n)+\tfrac12E-\log I_n
 =\frac{D_n^2s^4}{2\tau^2(s^2+\tau^2)^2},
\]
which is \(O(n^{-2})\) for bounded \(D_n\) and zero when \(D_n=0\).
The logarithm of the ratio for approximation (5) in that theorem is
\(\tfrac12\log(1+s^2/\tau^2)-s^2D_n^2/\{2\tau^2(s^2+\tau^2)\}\). This is
of exact order \(n^{-1}\) when \(D_n=0\).
\end{enumerate}
\end{hbtheorem}

The second item gives the one-step formula the same guaranteed order as
approximation (5), not a better one. The source says the formula often
improves on (5) \cite[Sec.~4.1.2, p.~778]{KassRaftery1995BayesFactors}.
The third item computes one instance of that improvement. It is not a
general comparison theorem. The formula uses the prior precision \(G\)
in place of the negative Hessian of the log prior. Those agree for a
normal prior, and the error bound does not require them to agree.

\section{The precise fixed-dimensional BIC consequence}
\label{sec:p15-bic-consequence}

\begin{hbtheorem}{thm:p15-bic-leading-order}{proved=annex:p15.bic-leading-order}{Log evidence and discarded constant-order terms}
Consider two models of fixed dimensions \(d_0,d_1\) satisfying the MLE
Laplace hypotheses with \(F_{k,n}=nJ_{k,n}\) and proper prior densities
\(\pi_k\). Define
\[
 c_{k,n}=\frac{d_k}{2}\log(2\pi)+\log\pi_k(\widehat\theta_{k,n})
                         -\frac12\log\det J_{k,n}.
\]
Then \(c_{k,n}=O(1)\) and
\[
 \log I_{k,n}
 =\ell_{k,n}(\widehat\theta_{k,n})-\frac{d_k}{2}\log n
                +c_{k,n}+O(n^{-1}).
\]
A parameter-free model with positive finite likelihood uses
\(I_{k,n}=L_{k,n}\), \(d_k=0,c_{k,n}=0\).
With
\(\mathrm{BIC}_k=-2\ell_{k,n}(\widehat\theta_{k,n})+d_k\log n\),
\[
 \log B_{10}
 =-\tfrac12(\mathrm{BIC}_1-\mathrm{BIC}_0)
                     +c_{1,n}-c_{0,n}+O(n^{-1}).
\]
Discarding the \(c\)-terms gives an \(O(1)\) error in log evidence,
not a relative \(O(n^{-1})\) approximation to the factor. If
\(c_{1,n}-c_{0,n}\to c_*\), the ratio of the BIC-based factor to
the exact factor tends to \(e^{-c_*}\). A relatively negligible error on
the log scale follows if \(|\log B_{10}|\to\infty\); it is not asserted
when the log factor remains bounded.
\end{hbtheorem}

Before discarding terms, the two-model MLE formula contains
\[
\begin{split}
2\log B_{10}={}&2(\ell_1(\widehat\theta_1)-\ell_0(\widehat\theta_0))
 +(d_1-d_0)\log(2\pi)-\log\det F_1+\log\det F_0\\
 &+2\log\pi_1(\widehat\theta_1)-2\log\pi_0(\widehat\theta_0)+O(n^{-1}).
\end{split}
\]
In particular both prior-log terms have coefficient two. The information
growth \(F_n=nJ_n\), rather than a convenient count of rows in a table,
determines the role of \(n\). Different blockwise information rates, singular
curvature, boundary modes and increasing dimension require other results
\cite[Sec.~4.1.3]{KassRaftery1995BayesFactors}.

The \(c\)-terms can be controlled by the choice of prior. Kass and
Raftery describe a prior on \(\psi\) of the form
\(\operatorname{normal}(\psi_0,W)\) with \(|W|=|\mathbf J(\psi_0)|^{-1}\),
where \(\mathbf J\) is the Fisher information. This prior carries the
information of one observation. Under null orthogonality and local
alternatives, the Schwarz criterion then approximates the log factor to
order \(O(n^{-1/2})\) \cite[Sec.~4.1.3, p.~779]{KassRaftery1995BayesFactors}.
Null orthogonality means that the Fisher information for
\((\beta,\psi)\) is block diagonal when \(\psi=\psi_0\). Local
alternatives means that \(\psi-\psi_0=O(n^{-1/2})\)
\cite[Sec.~5.2, p.~782]{KassRaftery1995BayesFactors}. The next theorem
states deterministic versions of these premises, evaluated along the
observed data sequence, and derives the order from the \(c\)-terms.

\begin{hbtheorem}{thm:p15-w02b-p15-056-unit-information-schwarz}{proved=annex:p15-w02b-p15-056-unitinfo}{The unit-information prior and an \(O(n^{-1/2})\) Schwarz approximation}
Let \(H_1\) have parameter \(\theta_1=(\beta,\psi)\in\mathbb R^{d_0}\times\mathbb R^p\),
\(p\geq1\), and let \(H_0\) be the submodel \(\psi=\psi_0\) with
parameter \(\beta\). When \(d_0=0\), \(\beta\) is absent and every
condition on it is vacuous. Suppose both models satisfy the hypotheses
of \cref{thm:p15-bic-leading-order}, with maximum-likelihood estimates
\((\widehat\beta_n,\widehat\psi_n)\) and \(\widehat\beta^*_n\) and
information rates \(J_{1,n},J_{0,n}\) whose eigenvalues lie in
\([c,C]\). Assume the following, with \(K\) independent of \(n\).
\begin{enumerate}
\item (Independence and a normal prior on \(\psi\).)
\(\pi_1(\beta,\psi)=\pi_0(\beta)\,\phi_W(\psi-\psi_0)\), where \(\phi_W\)
is the \(\operatorname{normal}(0,W)\) density, \(W\) is fixed and
positive definite, and \(\pi_0\) is the prior density under \(H_0\).
\item (Local alternative at the data scale.)
\(\|\widehat\psi_n-\psi_0\|\leq Kn^{-1/2}\) and
\(\|\widehat\beta_n-\widehat\beta^*_n\|\leq Kn^{-1/2}\), and
\(\log\pi_0\) is \(K\)-Lipschitz on the segment joining
\(\widehat\beta_n\) and \(\widehat\beta^*_n\).
\item (Null orthogonality and unit information.) In the block form
\(J_{1,n}=\left(\begin{smallmatrix}A_n&B_n\\B_n^T&D_n\end{smallmatrix}\right)\),
\(\|A_n-J_{0,n}\|\leq Kn^{-1/2}\), \(\|B_n\|\leq Kn^{-1/2}\) and
\(|\log\det D_n+\log\det W|\leq Kn^{-1/2}\).
\end{enumerate}
Then \(c_{1,n}-c_{0,n}=O(n^{-1/2})\), and for
\(S_n=\ell_1(\widehat\beta_n,\widehat\psi_n)-\ell_0(\widehat\beta^*_n)-\tfrac p2\log n
 =-\tfrac12(\mathrm{BIC}_1-\mathrm{BIC}_0)\),
\[
 \log B_{10}=S_n+O(n^{-1/2}).
\]
For the normal-mean instance \(d_0=0\), \(\psi_0=0\),
\(X_i\mid\psi\sim\operatorname{normal}(\psi,\sigma^2)\) with known
\(\sigma\), and \(W=\sigma^2\), write \(z=\sqrt n\,\bar x/\sigma\). Then
\[
 \log B_{10}-S_n=-\tfrac12\log(1+n^{-1})-\frac{z^2}{2(n+1)} .
\]
This is \(O(n^{-1})\) when \(|z|\leq K\). If instead \(\bar x\to\psi_1\neq0\),
the difference tends to \(-\psi_1^2/(2\sigma^2)\), so the local-alternative
premise cannot be dropped.
\end{hbtheorem}

The third premise uses the observed information block and the prior
covariance. The source's choice \(|W|=|\mathbf J(\psi_0)|^{-1}\) meets it
when the observed \(\psi\)-block is within \(O(n^{-1/2})\) of the Fisher
block in log determinant. As
\cref{thm:p15-w02b-p15-056-information-replacement-boundary} shows, that
comparison holds only along suitable data sequences.

\begin{hbtheorem}{thm:p15-laplace-normal-example}{proved=annex:p15.laplace-normal-example}{An exact normal calculation and three failure boundaries}
For \(X_n\mid\theta\sim\operatorname{normal}(\theta,\sigma^2/n)\),
\(\theta\sim\operatorname{normal}(\mu,\tau^2)\), with fixed positive
\(\sigma,\tau\), the posterior-mode Laplace calculation is exact.
Writing \(s^2=\sigma^2/n\), the MLE approximation is \(\pi(x_n)\) and
\[
 \frac{\widehat I_{\rm MLE}}{I_n}
 =\sqrt{1+s^2/\tau^2}\,
   \exp\!\left\{-\frac{s^2(x_n-\mu)^2}{2\tau^2(s^2+\tau^2)}\right\}.
\]
For a bounded sequence \(x_n\), the stated hypotheses hold for both
variants and this ratio is \(1+O(n^{-1})\).

For \(x_n=\mu=0,\sigma^2=1,\tau^2=4\), against the point-null model
\(\theta=0\), exact, posterior-Laplace, MLE-Laplace and BIC factors
\(B_{10}\) are respectively
\[
 (1+4n)^{-1/2},\quad(1+4n)^{-1/2},\quad(2\sqrt n)^{-1},\quad n^{-1/2}.
\]
The BIC-to-exact ratio tends to two.

The following proper-prior likelihood examples show why hypotheses matter.
For a uniform prior on \([0,1]\) and likelihood \(e^{-n\theta^2/2}\),
the boundary-mode integral is asymptotic to one half of
\(\sqrt{2\pi/n}\). For a uniform prior on
\((-2,-1/2)\cup(1/2,2)\) and likelihood
\(e^{-n(|\theta|-1)^2/2}\), the two-mode integral is asymptotic to
\((2/3)\sqrt{2\pi/n}\), twice one mode's contribution.
For a uniform prior on \([-1,1]\) and likelihood \(e^{-n\theta^4}\),
the integral is asymptotic to
\(n^{-1/4}\int_0^\infty e^{-u^4}du\), not a dimension-one
\(n^{-1/2}\) rate.
\end{hbtheorem}

These likelihoods can arise from \(n\) successes in Bernoulli experiments
whose success probabilities are the corresponding exponentials with \(n=1\).
They are not unnormalized parameter priors. In the examples, respectively,
the interior-ball, one-mode tail, and positive-curvature premises fail.

\section{Direct simulation and importance sampling}
\label{sec:p15-evidence-monte-carlo}

Fix observed data and suppress their symbol. Let \(\Pi\) be a proper prior,
\(L(\theta)\) the finite nonnegative full likelihood, and
\(I=\int L\,d\Pi\in(0,\infty)\). Put \(T(d\theta)=L(\theta)\Pi(d\theta)\)
and posterior \(P=T/I\). Simulation count is \(M\), not sample-information
index \(n\).

\begin{hbtheorem}{thm:p15-direct-evidence-mc}{proved=annex:p15.direct-evidence-mc}{Independent prior simulation}
For i.i.d. \(\theta_i\sim\Pi\),
\[
 \widehat I_M=M^{-1}\sum_{i=1}^M L(\theta_i)
\]
is unbiased and converges almost surely to \(I\).
If \(\int L^2d\Pi<\infty\), its variance is
\(\{\int L^2d\Pi-I^2\}/M\). Integrability of \(L\), rather than a
finite second moment, suffices for the strong-law conclusion.
\end{hbtheorem}

\begin{hbtheorem}{thm:p15-importance-evidence}{proved=annex:p15.importance-evidence}{Ordinary and self-normalized importance are different}
Let \(Q\) be a proposal probability and suppose \(T\ll Q\).
For i.i.d. \(\theta_i\sim Q\), choose a finite almost-everywhere version
\(V=dT/dQ\). Then \(M^{-1}\sum_iV(\theta_i)\) is unbiased, converges
almost surely to \(I\), and has variance
\(\{\int V^2dQ-I^2\}/M\) when that second moment is finite.

For the source's self-normalized form assume the stronger
\(\Pi\ll Q\), let \(W=d\Pi/dQ\), and use
\[
 \widetilde I_M=
 \frac{\sum_iW(\theta_i)L(\theta_i)}{\sum_iW(\theta_i)}
\]
when its denominator is positive. It is eventually defined almost surely
and converges to \(I\). It need not be unbiased for finite \(M\).
If both measures have densities, \(W=\pi/q\) on \(q>0\);
support domination is essential.
\end{hbtheorem}

Kass and Raftery's equation (10) is the self-normalized ratio. An ordinary
importance mean divides by \(M\), not by the random sum of weights.
A proposal improving concentration can worsen tail behavior; no claim of
improvement follows merely from calling a proposal an importance distribution
\cite[Sec.~4.2]{KassRaftery1995BayesFactors}.

\section{Posterior harmonic means and their support}
\label{sec:p15-harmonic-evidence}

\begin{hbtheorem}{thm:p15-harmonic-evidence}{proved=annex:p15.harmonic-evidence}{Reciprocal identities and independent-simulation convergence}
Let \(G\) be a probability measure with \(G\ll T\), and set
\(H=dG/dT\), with a finite version on the posterior support. Then
\[
 E_P H=1/I,\qquad
 E_P H^2=\frac1I\int H^2\,dT.
\]
For i.i.d. \(\theta_i\sim P\), the mean
\(A_M=M^{-1}\sum_iH(\theta_i)\) is unbiased for \(1/I\) and
converges to \(1/I\) almost surely. Its inverse, when positive, is
eventually defined and converges to \(I\); it is not generally unbiased.

The classical harmonic mean takes \(G=\Pi\). It requires
\(L>0\) \(\Pi\)-almost everywhere, gives \(H=1/L\), and has finite
reciprocal-weight variance exactly when \(\int L^{-1}d\Pi<\infty\).
If \(\Pi(L>0)=s<1\), the same posterior average of \(1/L\) instead
converges to \(s/I\), so its inverse converges to \(I/s\).
For density \(g\) of \(G\) and \(h=L\pi\) of \(T\),
the generalized weight is \(g/h\), and its second-moment requirement
is \(\int_{\{h>0\}}g^2/h<\infty\), with \(g=0\) almost everywhere
where \(h=0\).
\end{hbtheorem}

The average \(g/(L\pi)\) estimates the reciprocal evidence. Inverting
this average estimates evidence consistently under the stated sampling
assumptions, but inversion does not preserve unbiasedness. The distinction
is necessary when reading the printed equation (12) and its description
\cite[Sec.~4.3]{KassRaftery1995BayesFactors}.

\begin{hbtheorem}{thm:p15-evidence-mc-clt}{proved=annex:p15.evidence-mc-clt}{Normal limits with the necessary second moments}
For the independent sampling schemes above, the following limits hold,
where variance zero denotes point mass at zero.
An ordinary mean of weights \(Z\) with mean \(I\) and variance
\(v<\infty\) satisfies
\(\sqrt M(\widehat I_M-I)\Rightarrow\operatorname{normal}(0,v)\).
For self-normalized importance, if
\(v_{\rm SN}=E_Q[W^2(L-I)^2]<\infty\), then
\[
 \sqrt M(\widetilde I_M-I)\Rightarrow\operatorname{normal}(0,v_{\rm SN}).
\]
For a generalized harmonic mean, if \(v_H=\operatorname{Var}_P(H)<\infty\),
\[
 \sqrt M(A_M^{-1}-I)\Rightarrow\operatorname{normal}(0,I^4v_H).
\]
Assign any fixed finite default when a finite-sample denominator is zero;
those events disappear in probability under the stated strong laws.
No corresponding normal limit is asserted when these moment premises fail.
\end{hbtheorem}

\begin{hbtheorem}{thm:p15-mc-evidence-boundaries}{proved=annex:p15.mc-evidence-boundaries}{Calculated support, bias, variance and dependence failures}
The following examples use proper priors and finite likelihoods.
\begin{enumerate}
\item On two parameter values, let \(\Pi=(\varepsilon,1-\varepsilon)\),
\(0<\varepsilon<1\), and \(L=(1,0)\). Then \(I=\varepsilon\).
Direct simulation has relative variance
\((1-\varepsilon)/(M\varepsilon)\) and probability
\((1-\varepsilon)^M\) of returning zero. The posterior is concentrated
on the first value, so the classical harmonic estimate equals one,
not \(\varepsilon\). Proposal \(Q=(1,0)\) gives a valid zero-variance
ordinary importance estimate \(\varepsilon\), but the attempted
self-normalized estimate equals one.
\item With prior \((1/2,1/2)\), likelihood \((1,3)\), proposal
\((3/4,1/4)\), and \(M=1\), evidence is 2. Ordinary importance is
unbiased with variance \(16/3\); self-normalized importance has expectation
\(3/2\). The classical posterior harmonic estimate with one draw has
expectation \(5/2\), despite its reciprocal being unbiased for \(1/2\).
Likelihood values exceeding one are allowed for data densities.
\item Let \(\Pi=\operatorname{normal}(0,1)\), \(L=1\), and
\(Q=\operatorname{normal}(0,1/4)\). Ordinary importance has full support
and mean one but infinite variance.
\item For likelihood \(L(\theta)=\phi_{\theta,s^2}(x)\) and prior
\(\operatorname{normal}(0,\tau^2)\), the classical harmonic reciprocal
weight has infinite variance whenever \(\tau^2\geq s^2\), although
its expectation is finite and the i.i.d. strong law still applies.
\item For prior \((1/2,1/2)\), likelihood \((1/4,3/4)\), the evidence
is \(1/2\) and posterior is \((1/4,3/4)\). Initialize a chain in that
posterior and keep its state fixed forever. Every marginal is posterior,
but the harmonic estimate is forever \(1/4\) or \(3/4\), never \(1/2\).
\end{enumerate}
\end{hbtheorem}

All simulation convergence statements above use the integrable i.i.d.
strong law, not an unstated MCMC theorem. For dependent output one needs
a theorem giving the required ergodic averages, with its actual hypotheses.
Approximate posterior draws also need an approximation-error argument.
Stationarity alone is insufficient, as the last example proves.

\section{Mixture, bridge, Laplace--Metropolis and ratio estimators}
\label{sec:p15-w02b-p15-056-further-estimators}

The notation \(\Pi,L,I,T,P\) is that of
\cref{sec:p15-evidence-monte-carlo}. Kass and Raftery describe the
estimators below in words, and display only the Savage--Dickey ratio
\cite[Sec.~4.3, p.~780]{KassRaftery1995BayesFactors}. The definitions
here make those descriptions precise, and the theorems give the
conditions under which each estimator reaches its intended target.

\begin{hbdefinition}{def:p15-w02b-p15-056-mixture-estimators}{Prior/posterior mixture estimators}
Fix \(0<\delta<1\). For \(p>0\) put
\[
 \psi_p(\theta)=\frac{p-L(\theta)}{\delta p+(1-\delta)L(\theta)} .
\]
The importance function of the first estimator is the mixture
\(Q_\delta=\delta\Pi+(1-\delta)P\), whose density relative to \(\Pi\)
is \(\delta+(1-\delta)L/I\). Its weight in the self-normalized equation
(10) is \(w=1/\{\delta+(1-\delta)L/I\}\), and this weight contains the
unknown \(I\).
\begin{enumerate}
\item \(\widehat{\rm pr}_3\): draw \(\theta_1,\ldots,\theta_m\)
i.i.d. from \(Q_\delta\), and replace \(I\) in the weight by a candidate
value \(p\). With \(w_p=p/\{\delta p+(1-\delta)L\}\), equation (10)
then reads
\[
 p=\frac{\sum_iw_p(\theta_i)L(\theta_i)}{\sum_iw_p(\theta_i)} .
\] Since \(p>0\), this is equivalent
to \(\Psi_m(p)=\sum_{i=1}^m\psi_p(\theta_i)=0\). The estimator is the
root of this equation when the root exists and is unique.
\item \(\widehat{\rm pr}_4\): draw \(\theta_1,\ldots,\theta_m\) i.i.d.
from \(P\) only. Add \(\delta m/(1-\delta)\) imagined prior draws whose
likelihood is set equal to its prior expectation, taken as the candidate
\(p\). Each imagined draw has weight \(w_p=1\), and (10) becomes the
fixed-point equation
\[
 p=\frac{\delta mp/(1-\delta)+\sum_iw_p(\theta_i)L(\theta_i)}
         {\delta m/(1-\delta)+\sum_iw_p(\theta_i)} .
\]
Iterating the right side is the simple iterative scheme the source
mentions. The estimator is the fixed point when it exists and is unique.
\end{enumerate}
\end{hbdefinition}

\begin{hbtheorem}{thm:p15-w02b-p15-056-mixture-estimators}{proved=annex:p15-w02b-p15-056-mixture}{Existence, consistency and the limit of the mixture estimators}
\begin{enumerate}
\item For a sample of size \(m\), \(\Psi_m\) is continuous and
nondecreasing on \((0,\infty)\), with \(\Psi_m(\infty)=m/\delta\) and
\(\Psi_m(0+)=N_0/\delta-N_+/(1-\delta)\). Here \(N_+\) and \(N_0\) count
draws with positive and zero likelihood. A unique root exists exactly when
\((1-\delta)N_0<\delta N_+\). The fixed points of the
\(\widehat{\rm pr}_4\) equation are exactly the roots of
\(\sum_i\psi_p(\theta_i)=0\) over the posterior draws, and since
posterior draws have positive likelihood almost surely, that root exists
and is unique. In both cases bisection on the monotone function \(\Psi_m\)
computes the root.
\item \(\widehat{\rm pr}_3\) is eventually defined almost surely, converges
almost surely to \(I\), and satisfies
\(\sqrt m(\widehat{\rm pr}_3-I)\Rightarrow\operatorname{normal}(0,v_3)\),
where
\[
 v_3=\frac{E_{Q_\delta}\psi_I^2}
   {\bigl(E_{Q_\delta}[L/\{\delta I+(1-\delta)L\}^2]\bigr)^2}<\infty .
\]
The function \(\psi_p\) is bounded in absolute value by
\(\max\{1/\delta,1/(1-\delta)\}\), so no moment condition on \(L\) or
\(1/L\) is needed.
\item \(\widehat{\rm pr}_4\) converges almost surely to the unique root
\(p_\delta\) of \(E_P\psi_p=0\). It is consistent for \(I\) exactly when
\[
 \int\frac{L(I-L)}{\delta I+(1-\delta)L}\,d\Pi=0 .
\]
For the prior \((1/2,1/2)\) and likelihood \((1,3)\) on two points,
\(I=2\) but \(p_{1/2}=(1+\sqrt{13})/2\approx2.303\). For the same
example, \(\widehat{\rm pr}_3\) with \(\delta=1/2\) converges to \(2\).
\end{enumerate}
\end{hbtheorem}

In the second construction the imagined prior draws contribute nothing to
the equation \(\Psi_m(p)=0\). Their correct contribution would be
\(\delta m/(1-\delta)\) times \(E_\Pi\psi_p\), which is generally not
zero. At \(\delta=0\) the posterior equation is
\(\sum_i\{p/L(\theta_i)-1\}=0\), the harmonic-mean equation of
\cref{thm:p15-harmonic-evidence}. A small \(\delta\) keeps the limit
close to that equation but, as the example shows, a fixed positive
\(\delta\) need not give \(I\).

\begin{hbtheorem}{thm:p15-w02b-p15-056-bridge-identity}{proved=annex:p15-w02b-p15-056-bridge}{The basic bridge-type ratio identity}
Let \(\mu\) be a \(\sigma\)-finite measure on \(\Theta\). Let
\(q_1,q_2\geq0\) be measurable with
\(c_j=\int q_j\,d\mu\in(0,\infty)\) and probability densities
\(p_j=q_j/c_j\). Let \(\alpha\) be real and measurable, with
\(\int|\alpha|q_1q_2\,d\mu<\infty\) and
\(\kappa=\int\alpha q_1q_2\,d\mu\neq0\). Then
\[
 \frac{c_1}{c_2}=\frac{E_{p_2}[\alpha q_1]}{E_{p_1}[\alpha q_2]} .
\]
For independent i.i.d. samples \(\theta_i\sim p_1\), \(i\leq M_1\), and
\(\theta'_j\sim p_2\), \(j\leq M_2\), the ratio of the two sample means
of \(\alpha q_1(\theta'_j)\) and \(\alpha q_2(\theta_i)\) is eventually
defined and converges almost surely to \(c_1/c_2\) as
\(M_1,M_2\to\infty\). With \(q_1=L\pi\), so that \(c_1=I\), and a
probability density \(q_2=f\):
\begin{enumerate}
\item \(\alpha=1/f\) on \(\{f>0\}\) and zero elsewhere gives
\(E_f[L\pi/f]=I\cdot P(f>0)\), with \(P\) the posterior: the ordinary
importance identity with its support condition;
\item \(\alpha=1/(L\pi)\) on \(\{L\pi>0\}\) and zero elsewhere gives
\(1/E_P[f/(L\pi)]=I/\int_{\{L\pi>0\}}f\,d\mu\) whenever that integral
is positive: the generalized harmonic identity (12) with its support
condition.
\end{enumerate}
With \(q_1=L_1\pi_1\) and \(q_2=L_0\pi_0\) on a common parameter space,
the same identity gives \(B_{10}=c_1/c_2\) directly.
\end{hbtheorem}

Kass and Raftery note that Meng and Wong considered an optimal choice of
\(f\), computed iteratively, and that their framework applies to ratios
of integrals \cite[Sec.~4.3, p.~780]{KassRaftery1995BayesFactors}. The
identity above is the basic relation that such choices build on. The
optimal choice itself and its iteration are not derived in this chapter.

\begin{hbtheorem}{thm:p15-w02b-p15-056-laplace-metropolis}{proved=annex:p15-w02b-p15-056-lapmet}{When Laplace--Metropolis inherits Laplace accuracy}
Assume the posterior-mode construction of
\cref{thm:p15-laplace-mode-variants} satisfies
\cref{def:p15-laplace-family}, with \(q_n=\log(L_n\pi)=nh_n\), mode
\(\widetilde\theta_n\), \(H_n=nJ_n\), and the constants
\(\rho,c,C,A,c_0\). From simulation output, take a mode estimate
\(\theta^*_n\) with \(\varepsilon_n=\|\theta^*_n-\widetilde\theta_n\|<\rho\).
Take a covariance estimate
\(\Sigma^*_n=H_n^{-1/2}(\mathrm{I}+E_n)H_n^{-1/2}\) with \(E_n\)
symmetric and \(\|E_n\|\leq e_n<1/2\). Define the Laplace--Metropolis
value
\[
 \widehat I_{\rm LM}=(2\pi)^{d/2}|\Sigma^*_n|^{1/2}
 L_n(\theta^*_n)\pi(\theta^*_n).
\]
Then
\[
 \log\widehat I_{\rm post}-\log\widehat I_{\rm LM}
 =\{q_n(\widetilde\theta_n)-q_n(\theta^*_n)\}-\tfrac12\log\det(\mathrm{I}+E_n),
\]
with
\(nc_0\varepsilon_n^2\leq q_n(\widetilde\theta_n)-q_n(\theta^*_n)
 \leq\tfrac12n(C+A\rho)\varepsilon_n^2\) and
\(|\tfrac12\log\det(\mathrm{I}+E_n)|\leq de_n\). Hence
\(\log I_n-\log\widehat I_{\rm LM}=O(n^{-1})+O(n\varepsilon_n^2)+O(e_n)\).
The relative accuracy \(O(n^{-1})\) of \(\widehat I_{\rm post}\) is
inherited when \(\varepsilon_n=O(n^{-1})\) and \(e_n=O(n^{-1})\). If
\(e_n\to0\) and \(\liminf n\varepsilon_n^2>0\), the log error does not
tend to zero. A mode estimate with error comparable to the posterior
spread \(n^{-1/2}\) is therefore not enough.

In the normal model of \cref{thm:p15-laplace-normal-example}, take
\(\theta^*\) to be the maximizer of \(L_n\pi\) among \(M\) i.i.d. exact
posterior draws, and take \(\Sigma^*=H_n^{-1}\). Then
\(\log I_n-\log\widehat I_{\rm LM}=\tfrac12\min_{i\leq M}Z_i^2\) for
i.i.d. standard normal \(Z_i\). This error is positive almost surely, its
law depends on \(M\) but not on \(n\), and
\(\mathbb P(\tfrac12\min_iZ_i^2>t)=\{2-2\Phi(\sqrt{2t})\}^M\).
\end{hbtheorem}

The source's alternatives for the mode estimate are the posterior median
and nonparametric density estimation. For the covariance it recommends a
robust estimate \cite[Sec.~4.3, p.~780]{KassRaftery1995BayesFactors}.
All of these enter the theorem only through \(\varepsilon_n\) and
\(e_n\). Any accuracy claim for them needs bounds on those two quantities
as functions of both \(n\) and the simulation size.

\begin{hbtheorem}{thm:p15-w02b-p15-056-savage-dickey}{proved=annex:p15-w02b-p15-056-sdratio}{The Savage--Dickey density ratio and its prior condition}
Let \(H_1\) have parameter \((\beta,\psi)\in\mathbb R^{d_0}\times\mathbb R^p\),
prior density \(\pi_1(\beta,\psi)\), and a finite nonnegative jointly
measurable likelihood \(L(\beta,\psi)\). Let \(H_0\) be \(\psi=\psi_0\),
with prior density \(\pi_0\) for \(\beta\). Assume
\[
 I_1=\iint L\pi_1\in(0,\infty),\qquad
 I_0=\int L(\beta,\psi_0)\pi_0(\beta)\,d\beta\in(0,\infty).
\]
Evaluate densities at \(\psi_0\) using the fixed joint version:
\[
\begin{gathered}
 \pi_1(\psi_0)=\int\pi_1(\beta,\psi_0)\,d\beta\in(0,\infty),\qquad
 \pi_1(\beta\mid\psi_0)=\pi_1(\beta,\psi_0)/\pi_1(\psi_0),\\
 p_1(\psi_0\mid D)=I_1^{-1}\int L(\beta,\psi_0)\pi_1(\beta,\psi_0)\,d\beta .
\end{gathered}
\]
\begin{enumerate}
\item If \(\pi_0(\beta)=\pi_1(\beta\mid\psi_0)\) for almost every
\(\beta\), then \(p_1(\psi_0\mid D)>0\) and
\[
 B_{10}=\frac{\pi_1(\psi_0)}{p_1(\psi_0\mid D)} .
\]
\item In general, suppose \(L(\beta,\psi_0)\pi_1(\beta\mid\psi_0)=0\) for
almost every \(\beta\) with \(\pi_0(\beta)=0\). Then
\[
 \frac{p_1(\psi_0\mid D)}{\pi_1(\psi_0)}
 =B_{01}\,E_{P_0}\!\left[\frac{\pi_1(\beta\mid\psi_0)}{\pi_0(\beta)}\right],
\]
where \(P_0\) is the posterior under \(H_0\).
\item Replacing \(\pi_1(\cdot,\psi_0)\) by another nonnegative \(g\)
with \(\int g\in(0,\infty)\) changes \(\pi_1\) only on a Lebesgue-null
set. It leaves \(I_1\) and \(B_{10}\) unchanged, but changes the left
side of the second item to \(\int L(\beta,\psi_0)g(\beta)\,d\beta/(I_1\int g)\).
\item With \(d_0=0\), \(X\sim\operatorname{normal}(\psi,s^2)\) and
\(\psi\sim\operatorname{normal}(0,\tau^2)\) under \(H_1\), the first
item reproduces the factor of \cref{thm:p15-normal-point-null-factor}.
With \(X_1\mid\beta\sim\operatorname{normal}(\beta,1)\) and
\(X_2\mid\psi\sim\operatorname{normal}(\psi,1)\) independent, \(\beta,\psi\)
independent \(\operatorname{normal}(0,1)\) under \(H_1\),
\(\psi_0=0\) and \(\beta\sim\operatorname{normal}(0,4)\) under \(H_0\),
the condition fails. At \(x_1=0\) the density ratio equals
\(B_{01}\sqrt{5/2}\).
\end{enumerate}
\end{hbtheorem}

The source states the condition either as prior independence together with
its relation (13), or as \(\pi(\beta\mid H_0)=\pi(\beta\mid\psi=\psi_0,H_1)\)
\cite[Sec.~4.3, p.~780]{KassRaftery1995BayesFactors}. The third item
shows why both sides must use the same density versions. The source
itself records that the conditional-density prescription depends on the
parameterization, an instance of the Borel paradox
\cite[Sec.~5.1, p.~781]{KassRaftery1995BayesFactors}.

\begin{hbtheorem}{thm:p15-w02b-p15-056-latent-ratio}{proved=annex:p15-w02b-p15-056-latent}{The latent-data ratio estimator}
Fix the data and let \((\mathcal Z,\nu)\) be a \(\sigma\)-finite latent
space. For \(k=0,1\), let \(p_k(z)=\mathrm{pr}(D,z\mid H_k)\geq0\) be
measurable, with \(m_k=\int p_k\,d\nu=\mathrm{pr}(D\mid H_k)\in(0,\infty)\).
Let \(P^z_k(dz)=p_k\,d\nu/m_k\) be the posterior of \(z\) under \(H_k\),
and let \(B_{10}(z)=p_1(z)/p_0(z)\) on \(\{p_0>0\}\).
\begin{enumerate}
\item \(E_{P^z_0}B_{10}(z)=B_{10}\,P^z_1(p_0>0)\). This equals \(B_{10}\)
exactly when \(p_1=0\) \(\nu\)-almost everywhere on \(\{p_0=0\}\).
\item For i.i.d. \(z_i\sim P^z_0\), the average of the \(B_{10}(z_i)\)
converges almost surely to \(B_{10}P^z_1(p_0>0)\). Its variance is
finite exactly when \(\int_{\{p_0>0\}}p_1^2/p_0\,d\nu<\infty\), and then
\cref{thm:p15-clt-iid} gives its normal limit.
\item Let \(z=\theta_1\) be the parameter of \(H_1\), with prior density
\(\pi_1\), likelihood \(L_1\) and \(p_1=L_1\pi_1\). Let \(H_0\) be
parameter-free with likelihood value \(L_0>0\), and give \(z\) the
auxiliary law \(\pi_1\) under \(H_0\), so that \(p_0=L_0\pi_1\). With the
roles of the models exchanged, posterior draws \(z_i\) under \(H_1\) give
the average of \(L_0/L_1(z_i)\). This is \(L_0\) times the harmonic-mean
statistic of \cref{thm:p15-harmonic-evidence}, and it converges to
\(B_{01}\) exactly when \(L_1>0\) \(\pi_1\)-almost everywhere.
\end{enumerate}
\end{hbtheorem}

The third item is one precise form of the source's remark that, when
\(z\) is present in \(H_1\) but not in \(H_0\), the latent-data method
reduces to the harmonic-mean estimator of \(\mathrm{pr}(D\mid H_1)\)
\cite[Sec.~4.3, p.~780]{KassRaftery1995BayesFactors}. The source
simulates the \(z^{(i)}\) by Markov chain methods. The convergence
statements here assume independent draws, and the limitation stated
below applies.

Three developments described in Section~4 of the source are not
established here. The first is error bounds for adaptive Gaussian
quadrature, which the source describes as effective for peaked integrals
of modest dimension, roughly less than nine
\cite[Sec.~4.2, p.~779]{KassRaftery1995BayesFactors}. The second is the
optimal bridge function. The third is ergodic averaging for Gibbs,
Metropolis or weighted-likelihood-bootstrap output. Every simulation
limit above assumes independent draws. For dependent output, the
stationary identity chain of \cref{thm:p15-mc-evidence-boundaries} shows
that correct marginals alone do not yield these limits.
